% ============================================================
%  R. Nielsen, "Om theoretisk og praktisk Erkjendelse, nogle Bemærkninger"
%  [On Theoretical and Practical Cognition: Some Remarks]
%  Nordisk Universitets-Tidsskrift, Tillægshefte til tiende Aargang,
%  ed. A. Ingerslev. Copenhagen 1866 (printed by S. Trier), pp. 63–77.
%  TRANSLATION (English)
%  Source of truth: transcription.tex (Danish), COMPLETE.
%  Mirrors transcription.tex 1:1; every printed page carries \opage{N}
%  at the corresponding point in the English. See TRANSLATION-PLAYBOOK.md.
%  Translated 22 Sep 2026 in one batch.
%  The original sets Latin and German in roman (it has no italic); here
%  they are italicized per the playbook. Greek (Rom. 1:19–20, p. 74) is
%  copied verbatim from the transcription, printer's readings included.
%  Terms: Erkjendelse = cognition; Viden = knowledge; Videnskab = science;
%  Tro = faith; Villie = will; Vilkaarlighed = arbitrariness;
%  Oprindelighed = originality; Fordobling = duplication.
% ============================================================
\documentclass[12pt,a4paper]{article}
\usepackage[utf8]{inputenc}
\usepackage[T1]{fontenc}
\usepackage{libertinus}
\usepackage{libertinust1math}
\usepackage{textalpha}   % Greek typed directly (Rom. 1:19--20 on p. 74)
\usepackage[english]{babel}
\usepackage[protrusion=true,expansion=true]{microtype}
\usepackage{setspace}
\usepackage[margin=1.3in]{geometry}
\usepackage{fancyhdr}
\usepackage{marginnote}
\newcommand{\opage}[1]{\marginnote{\footnotesize\textit{[#1]}}}
\usepackage{hyperref}

\hypersetup{
  pdftitle={On Theoretical and Practical Cognition: Some Remarks},
  pdfauthor={R. Nielsen},
  hidelinks
}

\setlength{\headheight}{15pt}
\pagestyle{fancy}
\fancyhf{}
\rhead{Nielsen, \emph{On Theoretical and Practical Cognition} (1866)}
\cfoot{\thepage}
\renewcommand{\thefootnote}{*)}

\setstretch{1.3}

\title{\textbf{On Theoretical and Practical Cognition:\\[4pt]
       Some Remarks}}
\author{R.\ Nielsen\\[4pt]
  \small Translated from the Danish}
\date{\emph{Nordisk Universitets-Tidsskrift}, Supplement to vol.~10\\
  (Copenhagen, 1866), pp.\ 63--77}

\begin{document}
\maketitle
\thispagestyle{fancy}

\bigskip

% --- p. 63 ---
\opage{63}Religion and science are great powers in the world of
spirit. That these ideal powers have often and vehemently made war on
one another cannot surprise us when we consider that in origin and
essence they are so heterogeneous, and that each of them for itself,
when it comes to directing the education and cultivation of the human
race, lays claim to hegemony. Religion presupposes a revelation,
revelation a miracle; but science denies the miracle, rejects
revelation, and thereby comes to offend religion. How ruinous the
conflict is, and how beneficial it would be, not only for the life of
the congregation but for the life of the people and for the
development of culture as a whole, if the misunderstanding---the old
misunderstanding which, by setting the two spiritual powers against
each other in hostility, has through centuries helped to split and
weaken the noblest forces and to confuse the conception of that which
is and must be the final goal of all endeavors---could at last be
removed, everyone who has seriously thought about the matter will
surely grant.

Should it now really be possible that a misunderstanding so deeply
ingrown as that between faith and knowledge,
% --- p. 64 ---
\opage{64}religion and science, could ever be removed from the ground
up?

It seems almost impossible, when one properly considers how much is
required here for a good understanding. For a good understanding
between faith and knowledge it is not enough, namely, that the parties
at last bring themselves to tolerate each other, so that some in
society are allowed to cultivate science without troubling themselves
about the life of faith, while the great majority, without insight
into the laws of nature and of science, ingenuously accepts whatever
theologians and priests, as time and occasion suggest, find it
serviceable to explain to the faithful about the miracles of faith.

Such an agreement brings with it only a characterless intermediate
state. But a characterless intermediate state has something sadly
dulling for the life of spirit; the salutary thing is that under an
armed neutrality a spark can smolder which, before one knows it,
bursts into flame.

A good understanding between religion and science requires that the
two powers be able essentially to acknowledge each other and, in
mutual acknowledgment, agree to work together toward one and the same
highest goal. But for this to happen, the two principles, faith and
knowledge, must allow themselves to be \emph{united in one
consciousness}. The believer must then, without harm to his faith, be
able to become a knower, and the knower, without denying or halving
his knowledge, be a believer; for the all-round cultivated person it
must, in other words, be possible to give to faith what is faith's,
and to knowledge what is knowledge's.

Once the demand is plainly, freely and openly stated, the objections
will press forward with spiritual natural necessity. ``Should the
absolutely heterogeneous principles of faith and knowledge really
allow themselves to be united in one consciousness, then the same
person must indeed accept as believing what he as knowing denies, the
miracle of revelation; he must therefore at one and the same time both
acknowledge and not acknowledge the very same thing; but that is an
intolerable
% --- p. 65 ---
\opage{65}contradiction. Whoever, despite the contradiction, would
attempt to solve such a task must then begin by halving his
consciousness. Just as that king who smiled with half his face and
threatened with the other half showed himself, in a certain sense, as
a man with two faces, so the consciousness that was to unite the two
heterogeneous principles would have to be knowing with the one half
and believing with the other, have one truth in the one and another
in the other, and so show itself with two truths (\textit{duæ
veritates}). And what good would that do? So long as the two truths
are united in one consciousness only in such a way that they are kept
apart, each in its own room, by a kind of intellectual partition wall,
the opposition that earlier aroused strife and discord in the outward
has merely been moved into the inward. If the partition wall
bursts---and burst it surely must, once the spirit shakes off the
mechanical---then there is danger that the heterogeneous truths, as
they are united, will fuse together in an explosion that shatters
understanding and brain.''

By thus laying bare the supposed contradiction, we have managed to
see the problem come into being. But that a problem must come into
being is indeed the first condition for its being solvable. Can it,
then, be solved? It calls for an attempt.

The ``great, gruesome, terrible contradiction'' that people have
wished to find in the circumstance that consciousness, by
acknowledging two absolutely heterogeneous principles, would split
itself in such a way that it lost its unity and with it its energy,
does not, when we look more closely, amount to much. Anyone who has
practiced enough dialectic to have learned to see what it is that
essentially matters in the relation between an absolute and a
relative will without difficulty convince himself that principles
which are absolutely heterogeneous may very well stand in mutual
relation and in this relation show themselves relative. Anyone who
has acquired enough
% --- p. 66 ---
\opage{66}psychology to know the first basic features of the doctrine
of the unity of consciousness through its pervasive self-duplications
will, on some reflection, be able to tell himself that in the life of
spirit there are oppositions---deep, strong oppositions bordering on
dualism---which are nevertheless so far from weakening the unity of
self-consciousness that they on the contrary heighten and confirm it.

Among the commonest and therefore best-known forms of duplication of
the life of spirit is without doubt the opposition between theory and
practice, between theoretical and practical cognition.

That theory and practice must be able to stand in relation, indeed in
essential relation, to each other, no one will surely deny; but
notwithstanding this they each have their own principle and are to
that extent heterogeneous. If a theoretical grounding may not mix in
practical considerations, then neither can actual practice simply be
derived from theory. Were someone to ask: which is the original in the
life of consciousness, theoretical or practical cognition?---who would
then hesitate over an answer?

How pervasive the opposition between theoretical and practical
cognition is can best be seen from the corresponding transformation
of consciousness. The theoretical consciousness is impersonal, the
practical is personal; the former is contemplative, disinterested,
lost without self-regard in the universal, moved by necessity,
forgetting itself in the grounding of the objective; the latter, by
contrast, is interested, in conflict with existence, moved by fear
and hope, desire and expectation, full of self-regard even in its
self-denial. The theoretical consciousness is essentially
\emph{thinking}, the practical essentially \emph{willing}.

That will is required to develop thought, and thought to apply the
will, is understandable; thought and will, the theoretical and
practical factors of consciousness, certainly stand in mutual relation
and
% --- p. 67 ---
\opage{67}are, in their difference, to that extent relative; but
thought is not for that reason will, nor will thought; theory is not
practice, nor practice theory: the theoretical consciousness, which
gives itself up to thoughts and thinking, lives in another region than
the practical, which concentrates itself in resolve and will.

It is a misunderstanding---a misunderstanding deadly to religion,
morality and all higher development of character---to make the will a
modification of thinking, a misunderstanding which, as far as modern
philosophy is concerned, can be said to begin with Spinoza
(\textit{voluntas et intellectus unum et idem sunt}), to be completed
by Hegel. In the Hegelian philosophy it is \textit{das Denken und das
Denken und das allgemeine Denken} on which everything, first and last,
depends. From this standpoint J.\ L.\ Heiberg too, in his
``Introductory Lecture to the Logic Course,''\footnote{% printed mark *)
\emph{Prosaiske Skrifter}, vol.~I, p.~465.}
has set forth with consistency and clarity the relation between the
theoretical and the practical, the logical and the ethical. ``Only the
cultivation of thinking,'' he says, ``is the cultivation of the human
being, for thought is at once the \emph{human}, which separates us
from subordinate nature, and the \emph{divine}, which unites us with
the world above us. Perhaps it will be objected that the same can be
said of \emph{feeling}; so it can, but only insofar as feeling is
permeated by thought. Faith, love, honor, duty are certainly feelings
which separate us from nature and incorporate us into a higher world;
but what gives these feelings so great a value are the
thought-determinations contained in them, without which they could
not exist in the form of feelings. If, for example, there were not a
moral world in which good and evil were separated according to the
rational determinations of thought, then neither could we separate
them in our feeling; we would then have no
% --- p. 68 ---
\opage{68}moral feeling or feeling of duty, no love of the good and no
abhorrence of evil, for all these feelings presuppose a power that
determines them; this power is thought, which is the sole absolute
ground of determination.\,.\,.
% three close-set points, as printed
If it stands firm, then, that thought is both the human and the
divine in us, then it is equally certain that the cultivation of our
thinking is the cultivation of ourselves.''
When thinking is in this way made one with the human being, so that
the distinctive essence of the human being becomes only a
modification of universal, objective thinking, then it follows of
itself that the will, despite its subjective essence, must also be a
modification of objective thought.

But who does not see that such views must lead us onto extremely
dangerous byways? As opium acts as a narcotic on the nerves, so
objective thinking, when it is to be applied according to the above
theory, acts as a narcotic on the practical consciousness; and just as
the enjoyment of opium can at first replace the lost reality with
indescribable raptures and paradisal dream-visions but, when the dream
is over, leaves terrible effects behind, so too a deification of
theory's objective thinking can at first enrapture the beholder who
sees the ideals of existence in the mirror; but when the contemplation
is at an end, actuality stands only all the more angular, hard and
defiant against the right of ideality; it is then in vain that
thinking tries to explain away a reality it cannot bear.

In relation to actuality and the tasks of actual life, practical
cognition must absolutely have greater originality than theoretical.
The will, to be sure, in order to be will, must originally presuppose
thought; but, mark well, it is in the first place not theoretical but
practical thinking that the will presupposes, and in the second place
the will, with this presupposing of its own, is just as original as
thought. Just as in thought there is something that cannot be derived
from
% --- p. 69 ---
\opage{69}the will, since thought as thought is not will, so too in
the will there is something that cannot be derived from thought, since
will as will is not thought. Certainly such categories as freedom of
choice, choice, resolution, etc., fall under the illumination of
thought insofar as they become clear in moral cognition; but in
principle they are nonetheless categories of the will; their cognition
is essentially practical. What is cognized theoretically cannot be
cognized practically; it is no use appealing to the fact that the
content cognized is the same, only the mode of cognition different: in
every decisive cognition the content is determined by the mode, and
the mode by the content. That courage as a virtue can be cognized in
theory must be granted; but if the courage that is cognized
theoretically were really in each and every respect the same, the very
same, as that which is cognized practically, what would then follow?
It would indeed follow that an able theoretical thinker, without
venturing anything, without exposing his life to any real danger,
could think himself so far into true courage that in reality he
thereby thought true courage into himself, and thus, by objective
thinking, however cowardly he might otherwise be by nature, actually
became courageous. But that he who is in reality cowardly, however
theoretically he may think, can no more speculate himself into real
courage than he who is without money, even if he were to think
theoretically-objectively until he burst, can speculate himself into
real money, needs no proof.

How deep the chasm is which, because of the heterogeneous principles,
forms the separation between theoretical and practical cognition can
further be seen from the fact that absolute truth itself must divide
into two truths, an impersonal one, which corresponds one-sidedly to
thinking, and a personal one, which corresponds to the will. That the
idea of thinking is specially named truth, while the idea of the
will, according to a commonly accepted
% --- p. 70 ---
\opage{70}usage, is named the good, has in recent times not been
sufficient to avert the theoretical encroachment by which the two
ideas have been fused together in such a way that the good has
perished in the true, just as the will has in thinking. Theory does
indeed posit a difference between thought and will, between the true
and the good; but since the difference itself falls within theory, and
consequently is and remains theoretical, thinking naturally keeps the
upper hand throughout.

Every time the question is raised of the validity of the
presuppositions either of ethics or of religion, one notices the
presumption of theory. Freedom of choice, for example, cannot have
unconditional validity; and why not? Because objective thinking, when
it is to comprehend freedom, can do nothing other than reduce
arbitrariness to a moment subordinate to the preponderant motives, so
that freedom is transformed into a necessity transparent in
reflection. But the reality of the will surely demands an
unconditional self-determination; the idea of the good demands an
original freedom, valid in and through itself. Yes, answers theory,
let the will as factual arbitrariness make as merry as it likes, let
the idea of the good in practice pitch its demands as high as it
likes: when the question concerns the reality of the will and of the
good, then it is a matter of truth; and when it is a matter of truth,
then it is after all thinking---universally valid thinking, developing
itself in knowledge---that is and remains the supreme judge. If it is
the God of revelation, the almighty Creator of heaven and earth, that
is under discussion, theory knows at once that the ``almighty
Creator'' of the article of faith cannot possibly be the true God; and
why not? Because the act of creation cannot possibly be thought as an
act of arbitrariness, and so not, as the article nevertheless evidently
presupposes, as an unconditional act of freedom! But religion,
revealed religion, which proclaims itself as the absolute good, simply
cannot subsist without such a presupposition;
% --- p. 71 ---
\opage{71}God himself would indeed cease to be the Holy One, the truly
Almighty, if in his relation to the world he were not the absolutely
free God, the God who creates out of nothing. Religion here, religion
there! what good does it do, says theory, to appeal either to the
immediate authority of a revealed word or to the popular
representations of religion, when the question concerns an eternal
truth. Only for objective thought is truth objective; a god may be as
holy as he likes in the popular imagination; if he is not the God of
the concept, then neither is he a true God; if he is not a true God,
then his holiness means nothing.

Against such a despotism of knowledge a protest must be made. To
protest against the despotism of theories is by no means---as has
sometimes been alleged from an abstract theoretical standpoint---the
same as violating, or at least veiling, the rich truth; on the
contrary, it is high time, not only in the name of religion and
ethics but in the name of truth, indeed in the name of theory itself,
that objections be raised against the undue encroachment of objective
knowledge, of theoretical cognition.

The protest with which practical cognition repels the encroachment of
theory is not the desperate attempt of an intellectual tendency which,
for lack of objective grounds, would resort to a peremptory decree,
incite arbitrariness and appeal to a blind subjectivity. Practical
cognition is in this case so far from doing violence to theoretical
cognition that it begins, on the contrary, with freedom, the original
freedom which theory itself, when it properly comes to theoretical
reflection, must grant to subjectivity. For what is most subjective in
subjectivity is after all the self; but the selfhood of the self is
self-determining, is will. If theoretical cognition now finds itself
compelled to grant will and freedom the same originality it grants to
thinking and reflection, then it also finds itself compelled to
% --- p. 72 ---
\opage{72}grant the subjective, personal truth---the idea of the will,
of the good---the same originality it grants to the idea of thinking,
the impersonal, objective truth. ``And what more,'' asks the
theorist, ``can practical cognition then demand?'' It can demand, we
answer, that theory take its own concessions seriously; it can demand
that the practical idea, in everything practical, be just as
autonomous, and hence just as independent of irrelevant theories, as
the theoretical idea, in everything that exclusively concerns the
theoretical, must be exempt from the disturbing interference of
practical interests; it can demand that the principle of the will be
respected; it can demand that the coin of freedom---in spite of
objective, even dogmatic-objective, speculations---be converted from
theory's nominal value into the currency of character, into the
silver value of existence.

On these conditions, but also only on these, can a true understanding
be brought about between theory and practice. Once this understanding
has been brought about, not superficially but from the ground up, a
reconciliation of faith and knowledge, religion and science, will
follow as of itself; for religion is practical from first to last,
science theoretical from first to last: faith is an ideal of the will,
knowledge an ideal of thought.

It goes without saying that the ideals of thought and will must be
able to enter into just as many-sided an interrelation as thought and
will themselves; every attempt to draw a definitively fixed boundary
line that would separate theoretical cognition from practical is to
that extent unreasonable and meaningless. But there is here
nonetheless a boundary, a difference in principle, which arises from
the fact that in theory all the powers of the soul are subordinated to
thinking, whereas in practice they are subordinated to the will.

Of every idea of life, to be sure, there can be both a theoretical and
a practical cognition; but how
% --- p. 73 ---
\opage{73}pervasive the difference is becomes evident at a glance at
the idea of God. It is necessary to say it resolutely: the God of
thought, whose cognition is theoretical, cannot be the God of faith,
for of the God of faith there is only a practical cognition. The
opposition between faith and knowledge, religion and science,
practical and theoretical cognition, cannot be expressed more strongly
than when it is expressly said that they each have their own God.

``And this,'' one asks, ``is supposed to be the good understanding?''
We answer: Yes, it is the best understanding that can be conceived.
When science fully sees that it can no more cognize the God of faith
than religion can cognize the God of knowledge, while religion, from
the standpoint of its practical cognition, sees the same thing---what
then? ``Why, then,'' one exclaims, ``they must for the first time come
into serious conflict with each other, for each of them will naturally
take its God to be not merely a true God but the only true God!'' Such
effusions, which, as far as theory is concerned, characterize lively
rather than thorough heads, are evidently based on superficial
reasoning.

When it is said that faith and knowledge cannot possibly have the same
God, the meaning is naturally not that the God of faith is another
than the God of knowledge in the same sense as the Emperor of China is
another than the Emperor of Brazil; the meaning is rather that the
divine which is thought is not an object of faith, and that God,
insofar as he is an object of faith, is not, and by the nature of the
case cannot be, an object for comprehending thought.

That the science of spirit, insofar as it understands itself, must be
willing to enter into the proposed separation is natural; for no
philosophy has yet developed the concept of God in such a way that the
God cognized in the concept could become an object of faith. Whoever
claims that anyone can really believe either in
% --- p. 74 ---
\opage{74}Aristotle's God or in Leibnitz's God or in Hegel's God is
thoughtless and does not know what he is saying. But that religion,
for its part, cannot have anything against the separation either
likewise lies in the nature of the case, and is unequivocally granted
by the Apostle Paul.

When Paul (Rom.\ 1, 19--22) speaks of what is distinctive in the
heathens' cognition of God, he evidently lays weight on that in God's
being which can be thought, on ``that which may be known of God'' (τὸ
γνωστὸν τοῦ θεοῦ), on God's ``invisible being, his eternal power and
Godhead'' (ἣ τε ἀΐδιος αυτοῦ δύναμις καὶ θειότης), ``which are beheld
and understood'' (νοούμενα καθορᾶται)
% GREEK copied verbatim from transcription.tex, printer's readings
% (ἣ, αυτοῦ) included; see the note there.
from his works. It is therefore, in the stricter sense, not the God
of the will but the God of thought that has made himself known to the
heathens. It is the Godhead revealed in the realm of nature, the
Godhead of essentiality, of ground, of power, of laws and of rational
necessity, whose intellectual light has been able to shine for the
peoples of nature. But that a cognition of God which on these
conditions must become more theoretical than practical, and so aim
more at essence than at will, can furnish only a weak foundation for
the ethical---that the theory of God has not been able to secure the
principles of moral life---the following verses (22--32) make evident
in so penetrating a manner that the picture they sketch of corruption
must horrify us.

How different the heathen cognition of God is from the Jewish and the
Christian can be seen from the ninth chapter of the Epistle to the
Romans. For what kind of God is it, after all, of whom Paul speaks in
this chapter? It is not the God of thought, whose freedom dawns in an
overarching necessity; no, it is the God of freedom, whose essence is
unconditional will. The God ``revealed'' to the heathens has, to be
sure, a will too, but it is only a thought-will, a will that is a
moment of the good identical with objective truth (\textit{Deus vult,
quia bonum est}). The God of Israel, by contrast, is sheer will, his
thoughts will-thoughts,
% --- p. 75 ---
\opage{75}his conclusions decrees, his elections pure arbitrariness.
For that the God of revelation is represented as the God of
arbitrariness by no means implies that God's unconditional will should
play games either with necessity or with truth or with goodness; on
the contrary! All ground, all truth and goodness has precisely
(Rom.\ 11, 32--33) its center, its point of departure and its point of
rest in the divine will (\textit{bonum est, quia Deus vult}).

How impossible it is to develop any objective theory of the God of the
will, anyone who will look into the Epistle to the Romans can, on
some reflection, tell himself; and should he feel the need of further
proof, then the theological quackeries of scriptural interpretation,
the speculative self-contradictions of dogmatics and the abracadabra
of the theories of predestination will supply proofs in abundance. But
the same thing that makes it impossible for the God of arbitrariness
to be an object of theoretical cognition, that very same thing also
makes the revealed God of the will be, and necessarily be, an object
of faith. As object of faith the God of the will is originally the God
of miracle; for what is miracle but the unconditional power of the
divine will over all conditions? He who stakes his life on a general
thinking perhaps gets a theory about life, but is cheated of life
itself; he who sets his mind on the practical cognition of a holy will
perhaps misses the theory, but he gains life, the eternal life of
spirit, for the life of spirit is in the will.

How inwardly faith and practical cognition are bound up with the will
can be seen in obedience. To believe in God is to obey God, is to do
God's will; but to obey God is in turn the same as to believe in God,
for the first commandment of the will is faith. How inwardly faith and
obedience correspond to each other can be seen in the two most eminent
phenomena of religious cognition that our age has to show, Grundtvig
and S.\
% --- p. 76 ---
\opage{76}Kierkegaard: \emph{faith is essentially obedience}---this is,
in brief, Kierkegaard's explanation; \emph{obedience} is
\emph{essentially faith, faith in the Word of life that is the light
of men}---this is, in sum, Grundtvig's explanation. The two
explanations are so far from standing in mutual conflict that they
rather complete, illuminate and confirm one another.

Should it still really be, as some claim, impossible for one and the
same consciousness to acknowledge both the practical and the
theoretical, both faith and knowledge, then it would in any case be
the theoretical, not the practical, that would have to be given up.
Existence is practical; a consciousness can exist without being
theoretical; but a consciousness cannot be theoretical without
existing. It is therefore possible to put existence uppermost and keep
theory; but to derive existence from theory, to put theory uppermost
and keep existence, is impossible. The ideal of existence is will, for
existence is practical; the ideal of thought is knowledge, for
thought, objective thought, is theoretical: to put the ideal of
existence uppermost and keep the ideal of thought is possible, but to
put the ideal of thought uppermost and satisfy the ideal of existence
is impossible; in other words: to put faith uppermost and keep science
is possible; but to put science uppermost and keep faith is
impossible.

Should anyone, however, be so lively in the head and so light in his
thoughts that his head really spins, so that he cannot make out what
is to stand uppermost and what lowermost, faith or knowledge, but out
of sheer liveliness confuses uppermost with lowermost and lowermost
with uppermost, let him try to simplify matters by throwing away
faith, the ideal of the will, the principle of character, and living
on pure theory: he still does not escape the duplication. The
duplication of theoretical and practical is then placed within theory.
To preserve the unity of self-consciousness, a person would
% --- p. 77 ---
\opage{77}then have to have his theory in the practical, his practice
in the theoretical, his thinking in his willing and his willing in his
thinking; but no lively person will surely boast of having solved the
riddle of existence by solving such a task. By duplicating oneself
practically within the theoretical, one becomes, with all possible
liveliness, in lively duplication---not an ethical character, but an
intellectual amphibian.

\begin{center}\rule{0.2\linewidth}{0.4pt}\end{center}

\end{document}
